\section{Conclusions and limits}\label{sec:conclusions}

What can be concluded, and what cannot?

\paragraph{What we keep.}
\begin{itemize}
  \item The route, with the season, the hour and the wing level, explains \IsoExplained\% of the gap
        $\Delta\Hloc(1000~\mathrm{s})$ between group time and solo flight; the same for experienced
        pilots and beginners.
  \item What remains, \IsoDHloc{} \IsoDHlocBand{}, is small, with a band away from zero; at 100~s group
        time is less straight than the control.
  \item On the same route a group spends more time climbing and less gliding; the climb rate and the
        glide speed do not differ.
\end{itemize}

\paragraph{Limits.}
\begin{itemize}
  \item The days differ: season and hour are matched, the weather of the day is not.
  \item The controls flew with nobody within 10~km. They may differ from the pilots who fly in
        groups, and this test has no selection contrast to measure it.
  \item Only the \IsoPieces{} pieces whose route was flown alone on another day enter; their routes
        are frequented ones.
  \item The bands measure the sampling of pilots and blocks. The radius and the cover of the route
        were varied only in the general report, on a sample, and the isolation distance not at all.
\end{itemize}
Fixing the route removes most of the straighter progress of groups; whether what remains is caused by
the group stays open.
